\documentclass[a4paper]{article}
\usepackage[a4paper, landscape, margin=1cm]{geometry}
\usepackage{longtable}
\usepackage{array}
\usepackage[table]{xcolor}

\begin{document}
\pagestyle{empty}

\setlength{\LTleft}{0pt}
\setlength{\LTright}{0pt}
\setlength{\tabcolsep}{3pt}
\footnotesize
\begin{longtable}{|>{\raggedright\arraybackslash}p{0.08\textwidth}|>{\raggedright\arraybackslash}p{0.11\textwidth}|>{\raggedright\arraybackslash}p{0.11\textwidth}|>{\raggedright\arraybackslash}p{0.10\textwidth}|>{\raggedright\arraybackslash}p{0.11\textwidth}|>{\raggedright\arraybackslash}p{0.44\textwidth}|}
        \caption{Waste stream disposal methods and justifications by process area.}
        \label{tab:waste_stream_disposal_routes_landscape}                                                                                                                                                                                                                                                                                                                                                                                                                                                                                                                                                                                                                                                            \\
        \hline
        \textbf{Stream} & \textbf{Description}                                 & \textbf{Major components}            & \textbf{Total mass per batch}                                 & \textbf{Disposal method}                                                            & \textbf{Justification}                                                                                                                                                                                                                                                                                                                                                                                                                  \\
        \hline
        \endfirsthead
        \hline
        \textbf{Stream} & \textbf{Description}                                 & \textbf{Major components}            & \textbf{Total mass per batch}                                 & \textbf{Disposal method}                                                            & \textbf{Justification}                                                                                                                                                                                                                                                                                                                                                                                                                  \\
        \hline
        \endhead
        \hline
        \endfoot
        \hline
        \endlastfoot
        \rowcolor{gray!20}\multicolumn{6}{|c|}{\textbf{BATCH}}                                                                                                                                                                                                                                                                                                                                                                                                                                                                                                                                                                                                                                                        \\
        \hline
        F-24 and F-69   & Solvent waste from R-1 and R-2                       & \multicolumn{4}{l|}{Sent for DMF recycling}                                                                                                                                                                                                                                              \\
        \hline
        F-25 and F-72   & CO$_2$ waste from R-1 and R-2                        & CO$_2$                               & 40.4 kg                                                       & Vent to atmosphere                                                                  & CO$_2$ is non-toxic, non-flammable, and non-hazardous at the quantities generated, therefore it can be safely vented to the atmosphere via the plant ventilation system.                                                                                                                                                                                                                                                                \\
        \hline
        F-27 and F-74   & Spent resin beads from the fragment 1 AND 2 branches & DMF, resin beads                     & 207 kg DMF, 4.4 kg resin beads                                & Incineration                                                                        & Spent solid-phase synthesis resins (swollen with DMF) are contaminated with organic solvents and peptide residues, making them unsuitable for recycling or biological treatment. High-temperature incineration ensures complete destruction of organic contaminants and safe disposal.                                                                                                                                                  \\
        \hline
        F-31 and F-78   & Solvent waste from S-3 and S-7                       & Et$_2$O, TFA                         & 3742 kg Et$_2$O, 741.5 kg TFA                                 & Incineration + acid gas scrubber (NaOH)                                             & The stream contains large quantities of volatile organic solvents (Et$_2$O and TFA) which are unsuitable for wastewater discharge due to flammability, toxicity, and impact on marine life. Also, recovery by distillation would increase the costs of the plant marginally and pose additional hazards due to the flash point of Et$_2$O being -45 C. Hazardous waste incineration is therefore used to destroy the organic compounds. \\
        \hline
        F-38            & Solvent waste from S-4                               & Water, HCl, NaNO$_2$                 & 1005.7 kg water, 367 kg HCl, 4.6 kg NaNO$_2$                  & Oxidation + neutralisation, filtration, waste water disposal                        & The stream contains acidic aqueous waste (HCl) and nitrite (NaNO$_2$). Chemical oxidation using H$_2$O$_2$ converts NaNO$_2$ to NaNO$_3$ (6 kg), after which neutralisation using NaOH produces NaCl (590 kg), allowing the treated aqueous stream to meet conditions suitable for wastewater disposal.                                                                                                                                 \\
        \hline
        F-70            & Heavy metal and solvent waste from R-2               & DMF, DCM, Pd(PPh$_3$)$_4$ (Catalyst) & 2005 kg DMF, 1114 kg DCM, 19.7 kg Pd catalyst                 & Metal scavenging + filtration, incineration + acid gas scrubber (NaOH)              & Palladium catalyst residues are toxic and must be removed before disposal. Metal scavenging was done using 30 kg of carbon black followed by filtration recover the metal-carbon slurry, which was sent to a specialist for recovery (1.84 kg Pd). The remaining organic solvent mixture (DMF/DCM) is treated by hazardous waste incineration with acid gas scrubbing to neutralise acidic combustion gases.                            \\
        \hline
        \rowcolor{gray!20}\multicolumn{6}{|c|}{\textbf{CONTINUOUS}}                                                                                                                                                                                                                                                                                                                                                                                                                                                                                                                                                                                                                                                   \\
        \hline
        F-101 and F-120 & Solvent waste from S-9 and TFF respectively          & Water, desulfuriser solvents         & 4382 kg water, 66.6 kg desulfuriser solvents                  & Oxidation + neutralisation + biological treatment, filtration, waste water disposal & The stream is predominantly aqueous with low concentrations of organic compounds. Oxidation and neutralisation reduce the chemical oxygen demand (COD) and adjust pH, allowing the stream to be treated biologically, after which it is deemed safe to be discharged (sulphides converted to sulphates).                                                                                                                                \\
        \hline
        F-107 and F-108 & Solvent waste from RP-HPLC and freeze dryer          & Water, acetonitrile, organics        & 293 kg water, 134.5 kg acetonitrile, 20 kg organics           & Incineration                                                                        & The stream contains 4 wt\% acetonitrile, which exceeds concentrations typically tolerated in biological wastewater treatment and could inhibit microbial activity. Therefore, hazardous waste incineration is used for safe destruction of this stream.                                                                                                                                                                                 \\
        \hline
        \rowcolor{gray!20}\multicolumn{6}{|c|}{\textbf{SOLVENT RECOVERY}}                                                                                                                                                                                                                                                                                                                                                                                                                                                                                                                                                                                                                                             \\
        \hline
        F-113 and F-117 & Solvent waste from S-12 and bottom of S-14           & DMF, Fmocs, TBTU, organics           & 221.5 kg DMF, 244.8 kg Fmocs, 94.5 kg TBTU, 281.8 kg organics & Incineration + acid gas scrubber (NaOH)                                             & These streams contain complex organic synthesis reagents and peptide residues, which are not suitable for recovery or biological treatment. High-temperature incineration with acid gas scrubbing ensures complete destruction of organic compounds and neutralisation of acidic combustion gases.                                                                                                                                      \\
        \cline{1-5}
        F-114           & Solvent waste from top of S-13                       & DMF, piperidine, DIEA                & 295.5 kg DMF, 895.7 kg piperidine, 68.2 kg DIEA               & Incineration                                                                        &                                                                                                                                                                                                                                                                                                                                                                                                                                         \\
\end{longtable}
\normalsize

\end{document}
